% ═══════════════════════════════════════════════════════════════════════════
% Chapter 2: Background and Literature Review
% ═══════════════════════════════════════════════════════════════════════════

\chapter{Background and Literature Review} \label{chap:background}

\section{Velocity-Based Training: Theoretical Foundations}

\subsection{The Load-Velocity Relationship}

The theoretical foundation of VBT rests on the well-established load-velocity relationship first systematically characterized by Zatsiorsky and colleagues~\cite{zatsiorsky1995}. For compound barbell exercises, there exists a quasi-linear relationship between the external load and the maximum velocity achievable during the concentric phase:

\begin{equation}
    v = v_0 \left(1 - \frac{L}{1RM}\right)
\end{equation}

where $v$ is the mean concentric velocity, $v_0$ is the theoretical maximum velocity at zero load, $L$ is the absolute load, and $1RM$ is the one-repetition maximum.

Gonz{\'a}lez-Badillo and colleagues~\cite{gonzalez-badillo2010} demonstrated this relationship across multiple exercises with coefficients of determination $R^2 > 0.98$, establishing mean velocity as a reliable proxy for relative intensity. The slope and intercept of this relationship are exercise-specific and athlete-specific, requiring individual calibration for precise application.

\subsection{Velocity Zones and Training Adaptations}

Weakley et al.~\cite{weakley2021} established discrete velocity zones corresponding to distinct strength qualities:

\begin{table}[h]
    \centering
    \caption{Velocity zones and corresponding strength qualities}
    \label{tab:velocity_zones}
    \begin{tabular}{lcc}
        \toprule
        \textbf{Strength Quality} & \textbf{MCV (m/s)} & \textbf{Approx. \%1RM} \\
        \midrule
        Absolute Strength    & 0.15 -- 0.50 & 85 -- 100\% \\
        Acceleratives Strength & 0.50 -- 0.75 & 70 -- 85\% \\
        Strength-Speed       & 0.75 -- 1.00 & 55 -- 70\% \\
        Speed-Strength       & 1.00 -- 1.30 & 40 -- 55\% \\
        Starting Strength    & $> 1.30$      & $< 40\%$ \\
        \bottomrule
    \end{tabular}
\end{table}

The practical implication is that measurement error exceeding 0.05 m/s risks prescribing training outside the intended physiological adaptation zone. This accuracy requirement drives the engineering specifications for any VBT measurement system.

\subsection{Velocity Loss as Fatigue Metric}

S{\'a}nchez-Medina and Gonz{\'a}lez-Badillo~\cite{sanchez-medina2011} demonstrated that velocity loss within a set correlates with neuromuscular fatigue. Training to specific velocity loss thresholds (e.g., terminating at 20\% loss from first repetition) produces:

\begin{itemize}
    \item Superior hypertrophic adaptations compared to training to muscular failure
    \item Reduced central nervous system fatigue
    \item Preserved movement velocity across the microcycle
    \item Lower perceived exertion for equivalent training stimulus
\end{itemize}

This finding has transformed set prescription from arbitrary repetition targets to objective fatigue-based termination criteria.

\section{Inertial Measurement Units: Principles and Technologies}

\subsection{MEMS Accelerometer Operation}

Micro-Electro-Mechanical Systems (MEMS) accelerometers operate on the principle of capacitance change due to proof mass displacement. The ICM-42688-P employed in this system uses a differential capacitive structure with closed-force feedback~\cite{invenSense2020}:

\begin{equation}
    C_{diff} = C_+ - C_- = \frac{2\epsilon_0 A \cdot d}{d^2 - x^2} \approx \frac{2\epsilon_0 A \cdot x}{d^2}
\end{equation}

where $\epsilon_0$ is permittivity, $A$ is plate area, $d$ is nominal gap, and $x$ is displacement proportional to acceleration.

The device offers programmable full-scale ranges of $\pm$2, $\pm$4, $\pm$8, and $\pm$16 g with corresponding noise densities of 20, 40, 80, and 160 $\mu$g/$\sqrt{\text{Hz}}$. The $\pm$16 g range, while noisiest, is selected for weightlifting applications where peak accelerations during explosive lifts can exceed 8 g.

\subsection{MEMS Gyroscope Operation}

The gyroscope in the ICM-42688-P employs a vibrating structure design using Coriolis force detection. When the proof mass oscillates at frequency $\omega_s$ and experiences angular velocity $\Omega$, the Coriolis acceleration is:

\begin{equation}
    a_c = 2 \cdot (\vec{\Omega} \times \vec{v_s})
\end{equation}

This induces secondary oscillation proportional to angular velocity, detected via capacitive sensing. The device offers $\pm$125, $\pm$250, $\pm$500, $\pm$1000, and $\pm$2000 deg/s ranges with $\pm$2000 deg/s selected to accommodate rapid barbell rotations during Olympic lifting.

\subsection{Sensor Error Sources}

\subsubsection{Bias Instability}

Bias represents the output offset when no acceleration/rotation is applied. Bias instability manifests as:

\begin{equation}
    y(t) = a_{true}(t) + b_a(t) + n(t)
\end{equation}

where $b_a(t)$ is slowly drifting bias and $n(t)$ is white noise. For the ICM-42688-P, typical bias instability is 1 mg over temperature ranges of 0-60$^\circ$C, requiring periodic recalibration.

\subsubsection{Scale Factor Error}

Scale factor non-linearity causes measured values to deviate from true values at higher accelerations. The ICM-42688-P specifies $\pm$1\% scale factor accuracy, introducing systematic error that grows with acceleration magnitude.

\subsubsection{Temperature Dependence}

Both bias and scale factor exhibit temperature dependence. The internal temperature sensor enables compensation, but residual thermal drift of 0.5 mg/$^\circ$C remains after factory calibration.

\section{Attitude Estimation Algorithms}

\subsection{Quaternion Representation}

Quaternions provide singularity-free 3D rotation representation, preferred over Euler angles for sensor fusion applications. A unit quaternion $q$ is defined as:

\begin{equation}
    q = w + xi + yj + zk = [w, \vec{v}] = [\cos(\theta/2), \hat{n}\sin(\theta/2)]
\end{equation}

where $\theta$ is rotation angle and $\hat{n}$ is rotation axis. The ICM-42688-P outputs data in body frame; we seek the quaternion $q_{b\to w}$ mapping body-frame vectors to world frame.

\subsection{Madgwick Filter}

Madgwick et al.~\cite{madgwick2010} proposed a gradient descent approach minimizing the error between measured and predicted gravity/magnetic field directions. The quaternion derivative is:

\begin{equation}
    \dot{q} = \frac{1}{2} q \otimes \omega - \beta \cdot \frac{\partial f}{\partial q}
\end{equation}

where $\omega$ is gyroscope-derived angular velocity, $\beta$ is step-size parameter, and $\frac{\partial f}{\partial q}$ is gradient of constraint error. While computationally efficient (approximately 500 $\mu$s per iteration on Cortex-M4), the single $\beta$ parameter cannot dynamically adapt to motion vs. drift, introducing the -0.2 m/s peak velocity bias observed in early VBT pipelines~\cite{laidig2021}.

\subsection{Mahony Filter}

Mahony et al.~\cite{mahony2008} introduced proportional-integral feedback on the cross-product error between measured and estimated gravity direction. The PI structure explicitly estimates gyro bias as a state variable:

\begin{align}
    e &= \vec{a}_{meas} \times \vec{a}_{est} \\
    \vec{b}_{int} &= \vec{b}_{int} + K_i \cdot e \cdot dt \\
    \vec{\omega}_{corr} &= \vec{\omega}_{meas} - \vec{b}_{int} + K_p \cdot e
\end{align}

This structure separates bias estimation from attitude estimation, providing better performance under sustained motion.

\subsection{Error-State Kalman Filter (ESKF)}

The ESKF~\cite{barrau2015} decomposes the state into nominal and error components:

\begin{align}
    \text{Full state: } & X = [q, \vec{b}_g] \\
    \text{Error state: } & \delta X = [\delta\vec{\theta}, \delta\vec{b}_g]
\end{align}

The prediction step propagates the nominal state using gyro measurements while the error state follows linearized dynamics:

\begin{equation}
    \delta\vec{\theta}_{k+1} = \delta\vec{\theta}_k - \delta\vec{b}_{g,k} \cdot dt + \vec{w}_g
\end{equation}

The update step incorporates accelerometer observations when the acceleration magnitude indicates near-gravity conditions ($|\vec{a}| \in [0.9g, 1.1g]$). The Kalman gain optimally weights prediction vs. observation based on covariance matrices.

\subsection{VQF Filter}

Laidig et al.~\cite{laidig2021} developed the Versatile Quaternion-based Filter (VQF) specifically addressing VBT requirements. Key innovations include:

\begin{itemize}
    \item \textbf{Adaptive Accel Weight:} Trust factor $w_a$ varies based on $|\vec{a} - 1g|$, down-weighting accelerometer during high-dynamic motion.
    \item \textbf{Integrated Rest Detection:} Simultaneous monitoring of accel variance and gyro magnitude identifies rest intervals for bias estimation.
    \item \textbf{Decomposed Inclination/Yaw:} Separate processing of inclination (from accel) and yaw (from gyro integration) avoids the ``magnetic dead zone'' problem.
\end{itemize}

Benchmarking against 47 sessions demonstrated VQF achieves 0.06 m/s mean absolute error for peak concentric velocity vs. Madgwick's 0.35 m/s, a 5.8$\times$ improvement enabling reliable VBT zone classification.

\section{Zero-Velocity Update (ZUPT) Strategies}

\subsection{ZUPT Principle}

During rest intervals, the barbell velocity is known to be zero. This observation provides a powerful correction to accumulated integration drift. The measurement model is:

\begin{equation}
    \vec{z} = \vec{v}_{meas} + \vec{v}_{bias} + \vec{n}_v
\end{equation}

where during rest, $\vec{z}$ should equal zero, revealing velocity bias for correction.

\subsection{Rest Detection}

Reliable ZUPT requires accurate rest detection. Industry practice employs dual thresholds:

\begin{align}
    \sigma_a &< \theta_{acc} \quad \text{(acceleration variance)} \\
    \sigma_\omega &< \theta_{gyr} \quad \text{(gyro variance)} \\
    t_{hold} &> \theta_{time} \quad \text{(sustain duration)}
\end{align}

VQF employs adaptive thresholds based on motion intensity, with typical values of $\theta_{acc} = 0.02g$, $\theta_{gyr} = 2$ deg/s, and $t_{hold} = 80$ ms for VBT applications.

\subsection{Application to Velocity Integration}

Upon rest detection, the velocity estimate is corrected:

\begin{equation}
    \vec{v}_{corrected} = \vec{v}_{estimated} - \vec{k} \cdot \vec{v}_{estimated}
\end{equation}

where $\vec{k}$ is a gain matrix determining correction aggressiveness. Over-aggressive correction introduces jerk; overly conservative correction fails to arrest drift.

\section{Hardware Synchronization Techniques}

\subsection{Software Timestamping}

Software timestamping assigns arrival time at the application layer:

\begin{quote}
\textit{Pros:} Simple implementation, no additional hardware\\
\textit{Cons:} Jitter from OS scheduling (10-100 ms), network latency variance
\end{quote}

Inadequate for IMU-camera fusion requiring sub-millisecond alignment.

\subsection{Hardware Trigger Distribution}

Dedicated synchronization signals provide deterministic timing. The approach employed in this system uses the camera's FRAME\_SYNC output:

\begin{enumerate}
    \item Intel RealSense D455 asserts CMOS FRAME\_SYNC at exact exposure start
    \item NPN transistor converts 1.8 V logic to 3.3 V for ESP32 compatibility
    \item Concurrent distribution to ESP32 GPIO and IMU FSYNC pin
    \item IMU marks samples with timestamp relative to trigger edge
\end{enumerate}

Achieved alignment is $\pm$200 $\mu$s limited by IMU internal timestamp resolution.

\section{Wireless Telemetry Protocols}

\subsection{ESP-NOW Protocol}

ESP-NOW is a proprietary protocol developed by Espressif for peer-to-peer communication between ESP32 devices~\cite{espressif2023}:

\begin{itemize}
    \item \textbf{Connectionless:} No handshake overhead, 0.5 ms latency budget
    \item \textbf{Low Power:} Radio duty cycle <10\% for typical telemetry
    \item \textbf{Encryption:} AES-128 optional for data integrity
    \item \textbf{Multi-STA:} Single master can receive from 8+ slaves simultaneously
\end{itemize}

The proprietary nature precludes full protocol analysis, but impulse response measurements show 2.4 GHz Field-Test reliability up to 15 m through 2 walls.

\subsection{BLE Comparison}

Bluetooth Low Energy offers standardization but higher latency (5-15 ms connection interval minimum) and lower throughput (251 bytes/packet vs. ESP-NOW's 250 bytes with lower overhead).

\section{ASIC and FPGA Acceleration}

\subsection{Why Hardware Acceleration?}

This system targets mobile/edge deployment where power efficiency determines operational viability:

\begin{align}
    \text{Power}_{software} &\approx 1.5 \text{ W (ESP32 at full load)} \\
    \text{Power}_{ASIC} &\approx 0.185 \text{ W (estimated 28nm)} \\
    \text{Efficiency Gain} &\approx 8\times
\end{align}

For battery-operated operation, this translates to 3 hours vs. 24 hours continuous runtime.

\subsection{FPGA Implementation Strategy}

Xilinx Artix-7 provides cost-effective acceleration with sufficient logic density:

\begin{itemize}
    \item DSP slices: 240 (used for matrix-vector multiply in ESKF)
    \item Block RAM: 270 KB (covariance matrices, state buffers)
    \item LUTs: 172,800 (quaternion operations, filter logic)
\end{itemize}

Parallelism opportunities:
\begin{itemize}
    \item Concurrent quaternion rotation for multiple samples
    \item Pipeline Kalman filter update equations
    \item Simultaneous batch processing of rep windows
\end{itemize}

\subsection{ASIC Microarchitecture}

28nm CMOS process offers favorable power-performance tradeoff:

\begin{itemize}
    \item Operating frequency: 200 MHz (sufficient for 1 kHz IMU with 200$\times$ headroom)
    \item Logic layers: 9 metal layers for complex routing
    \item Power gating: Power down unused filter variants
    \item SRAM: 64 KB embedded for state storage
\end{itemize}

\section{Summary}

This chapter reviewed the theoretical foundations of VBT, MEMS sensor operating principles, attitude estimation algorithms, and hardware acceleration strategies. The VQF filter with dual-trigger rep segmentation emerges as the optimal algorithm choice for VBT, validated in Chapter 9. Chapters 10-11 extend this foundation to FPGA and ASIC implementations.

[Continue with ...]